% Unit 6 map -- one page.
\documentclass{apchem-worksheet}

\apcunit{6}
\apcunittitle{Thermochemistry}
\apcdoctitle{Unit Map}
\apcsubtitle{CED Unit 6 \textbullet\ 7--9\% of the AP Exam \textbullet\ 5 blocks + exam}

\begin{document}
\apctitleblock

\begin{apcnote}[How this unit is graded --- and what makes it different]
Unit exam \textbf{Fri Dec 11}: 20 multiple choice (\textbf{no calculator})
$+$ 2 short free-response, 48 minutes, 28 points.

This is the shortest unit in the course and the one with the cleanest
one-to-one match to the textbook --- Zumdahl Chapter 6 is essentially the
whole unit, with two topics borrowed from elsewhere.

What makes it different is that almost every question is a
\textbf{sign question}. Getting $\Delta H$ right is usually easy; getting
the sign right, and saying which way heat flowed and from whose point of
view, is where the points actually are.
\end{apcnote}

\section*{\sffamily Block calendar}

\renewcommand{\arraystretch}{1.3}
\noindent\begin{tabular}{@{}p{2.1cm}p{5.0cm}p{1.5cm}p{3.9cm}p{1.9cm}@{}}
\toprule
\textbf{Date} & \textbf{Topics} & \textbf{CED} & \textbf{Read (PDF pp.)} & \textbf{Practice}\\
\midrule
Mon Nov 30 & Endo/exothermic; energy diagrams; heat transfer & 6.1--6.3 & \S6.1 \textbullet\ 283--290 & WS 1\\
Wed Dec 2 & Heat capacity and calorimetry & 6.4 & \S6.2 \textbullet\ 290--297 & WS 2\\
Fri Dec 4 & Energy of phase changes; heating curves & 6.5 & \S10.8 \textbullet\ 523--530 & WS 2\\
Mon Dec 7 & Enthalpy of reaction; Hess's law & 6.6, 6.9 & \S6.2--6.3 \textbullet\ 290--301 & WS 3\\
Wed Dec 9 & Bond enthalpies; enthalpies of formation & 6.7, 6.8 & \S8.8 \textbullet\ 412--415; \S6.4 \textbullet\ 301--308 & WS 4 (FRQ)\\
\textbf{Fri Dec 11} & \textbf{UNIT 6 EXAM} & all & review sheet & ---\\
\bottomrule
\end{tabular}

{\sffamily\footnotesize Dates from the co-op calendar. Benchmark B and the
semester wrap follow in the week of Dec 14, before Christmas break.}

\vspace{0.5em}
\section*{\sffamily By the end of this unit you can\ldots}

\begin{itemize}[leftmargin=*,itemsep=0.12em]
  \item Classify a process as endothermic or exothermic from the sign of
        $q$ or from an observation. \ced{6.1}
  \item Draw and interpret an energy diagram for a reaction. \ced{6.2}
  \item Explain heat transfer and thermal equilibrium at the particle
        level. \ced{6.3}
  \item Use $q = mc\Delta T$ and calorimetry data to find $\Delta H$ per
        mole. \ced{6.4}
  \item Calculate the energy of a phase change and read a heating curve.
        \ced{6.5}
  \item Relate $\Delta H$ to the amount of substance reacting. \ced{6.6}
  \item Estimate $\Delta H$ from average bond enthalpies. \ced{6.7}
  \item Calculate $\Delta H^\circ$ from enthalpies of formation. \ced{6.8}
  \item Apply Hess's law to combine reactions. \ced{6.9}
\end{itemize}

\begin{apctrap}[Where students lose points in this unit]
Losing the sign of $q$ when switching between system and surroundings;
using $q = mc\Delta T$ \emph{during} a phase change, where the temperature
does not change; forgetting to divide by moles when reporting
\si{\kilo\joule\per\mole}; reversing a reaction in Hess's law without
flipping the sign; using \emph{formed minus broken} for bond enthalpies
(it is broken minus formed); and treating an element's
$\Delta H_f^\circ$ as anything other than zero.
\end{apctrap}

\end{document}
